\section{Related Work}
\paragraph{Leaderboards}
Our work studies how preference-based leaderboards coordinate firms’ incentives and, in turn, shape diversity.
Leaderboards crowdsource language-model preferences and provide granular information about preference~\citep{chatbot_arena,copilot_arena}.
Designing such leaderboards raises at least two challenges: aggregating heterogeneous preferences into a representative ranking~\citep{informative_leaderboards} and making that ranking robust to strategic manipulation. 
The first challenge is closely related to estimating top-$k$ rankings from limited pairwise-comparison data~\citep{spectral_mle,condorcet_pairwise}, and can be viewed as a social choice problem~\citep{pluralistic_leaderboards,evaluation_social_choice}.
For the second challenge, prior work has shown that leaderboards can be manipulated through vote rigging~\citep{vote_rigging}, strategic candidacy~\citep{leaderboard_illusion,strategic_candidacy}, and targeted manipulation~\citep{dropping_preferences}.
% TODO: Check the last sentences
Closest to our theoretical setting is \citet{leaderboard_incentives}, which studies how benchmark design induces strategic post-training and asks whether the resulting equilibrium ranks models according to their latent quality. 
We instead ask how heterogeneity in developers' production technologies interacts with a fixed common evaluation. 
We model costly investment through heterogeneous production functions, which allows us to characterize when production-side comparative advantage sustains specialization, and to study the dynamics of catch-up and frontier crowding.

\paragraph{Contest Theory}
Contest theory studies how strategic agents expend costly effort to improve their probability of success~\citep{borel_contest_theory,advertising_contest,contest_success_function}. 
Our model is closely related to multi-battle contests~\citep{tullock_contest,multi_battle_contests}, in which agents allocate costly resources across multiple competitive dimensions. 
A key difference is that we separate investment from competitive performance: developer-specific production functions \(F_i\) map multidimensional investment into capability profiles, which are then compared through the contest. 
This allows us to study how heterogeneity in production technologies interacts with a common evaluation signal to determine whether equilibrium outputs become similar or specialized, rather than focusing primarily on effort allocation or contest/prize design.

\paragraph{Strategic Machine Learning}
In many real-world settings, agents respond strategically to machine learning algorithms, as deploying prediction algorithms might either impact agent responses for future predictions themselves~\citep{strategic_classification,strategic_investment,performative_prediction,strategic_ranking}.
These strategic responses can also impose social costs, for example by inducing agents to expend additional effort simply to maintain favorable classifications~\citep{social_cost_strategic_classification}.
Our setting reverses the usual direction of strategic response: rather than individuals adapting to a deployed model, model developers adapt their investments to an evaluation mechanism.

\paragraph{Competition and Monoculture}
Finally, our work connects to a growing literature on algorithmic monoculture and the effects of competition on diversity.
Algorithmic monoculture describes settings in which multiple decision makers rely on similar algorithms, rankings, or predictions, causing their decisions to become correlated~\citep{correlated_errors,subjectivity_monoculture,monoculture_outcome}.
Prior work has shown that competition can act as an anti-coordination force, encouraging developers to differentiate~\citep{competition_diversity}.
We identify a countervailing force; when competitors are evaluated against a common benchmark, evaluation can coordinate investment toward the same capabilities.
Related work similarly shows that competition among increasingly capable models need not increase diversity and can, in some settings, reduce welfare~\citep{bayes_risk_welfare}.
Closest to our diversity question is \citet{aggregation_diversity}, which studies how aggregation and evaluation objectives affect producer incentives and consumer utility, showing that standard win rate can encourage homogenization while alternative weighted win-rate objectives can encourage specialization. 
We instead model the underlying investment process explicitly: developers allocate costly resources across capabilities through heterogeneous production functions. 
This allows us to characterize how common evaluation interacts with production-side comparative advantage to generate monoculture or specialization, and to study the dynamics of catch-up and frontier crowding.